% 03_concepts.tex - the two couplings and the unsaturated columns
\section{Two ways to couple the soil and the aquifer}
Read this section before you start. It explains what the steps change.
\begin{description}
\item[\cmd{CONCEPTUAL} (the default, as in Workshop~2)] The soil and the aquifer meet only in two ways: through recharge
  (whatever Raven puts into its \cmd{GROUNDWATER} store) and through a seepage drain at the land surface. Where the water
  table reaches the surface, MODFLOW sends groundwater back to Raven. The water table never enters the soil.
\item[\cmd{PHYSICAL}] Raven computes the water in the soil layers you list; MODFLOW starts below them. The top of the
  MODFLOW model is the base of these layers (the land surface minus their thickness). The water table can rise into the soil. There, recharge is rejected (it stays in
  the soil), saturated soil drains to the aquifer, groundwater rises into the soil, and a full soil zone overflows to the
  surface. Within each day, Raven and MODFLOW repeat their calculation until the soil and the aquifer agree. On sloping or
  uneven land, each cell samples the DEM at 5$\times$5 points, so the low points saturate first.
\end{description}
Below the soil zone, recharge can cross an unsaturated zone. \cmd{:SoilZoneVadose} adds a column to each cell, between the
base of the soil zone and a water table deeper than 5~cm. The column delays and spreads the recharge. It solves Richards'
equation (\cmd{RICHARDS}), a kinematic wave (\cmd{KINEMATIC}) or a cascade of linear stores (\cmd{LAG}).
